\subsection{6. Spurnetzwerk-Normen}\label{spurnetzwerk-normen}

\[\adjustbox{max width=\linewidth}{$\displaystyle 
\frac{\nu_n}{(n!)^2}=\sum_{|m|=n}\frac{\|Q^{\dagger m}F\|^2}{m!}
=\Big[\text{Koeffizient von }s^n\Big]\int d\mu(z)\,
\exp\Big(-\tfrac12\sum_{k\ge1}\frac{s^k}{k}\operatorname{tr}\big((\bar ZM)(ZM)\big)^k\Big),
$}\]

über bosonische Kohärenzzustände, die Onishi-Determinante det(1 −
M̄M)\^{}\{1/2\} und Wick-Paarung der z-Variablen. Jedes Netzwerk ist eine
Kontraktion von n Kopien des Tensors Φ\_\{ij,kl\} =
Σ\_A(M\_A)\emph{\{ij\}(M\_A)}\{kl\} (Einträge in \{−1,0,1\}). Netzwerke
werden als 4-Bein-Multigraphen unter Knotenumbenennung und Slot-Tausch
kanonisiert (n=4: 120 → 27 Klassen). Exaktheit: float64-Kontraktion mit
Absolutwert-Begleitkontraktion (\textbar Σa\_kb\_k\textbar{} ≤
Σ\textbar a\_k\textbar\textbar b\_k\textbar{} \textless{} 2⁵³ je
Schritt), sonst Auswertung modulo fünf Primzahlen \textless{} 2¹³ mit
CRT (Eindeutigkeit aus \textbar Netzwerk\textbar{} ≤
60ⁿ·64\^{}\{\#Zyklen\}). Validierung: ν₁, ν₂ per Brute Force, ν₃ gegen
die Vektorrechnung, CRT-Route allein reproduziert ν₃.

\subsection{7. T1-T8}\label{t1-t8}

{\def\LTcaptype{none} % do not increment counter
\begin{longtable}[]{@{}
  >{\raggedright\arraybackslash}p{(\linewidth - 4\tabcolsep) * \real{0.3333}}
  >{\raggedright\arraybackslash}p{(\linewidth - 4\tabcolsep) * \real{0.3333}}
  >{\raggedright\arraybackslash}p{(\linewidth - 4\tabcolsep) * \real{0.3333}}@{}}
\toprule\noalign{}
\begin{minipage}[b]{\linewidth}\raggedright
Tor
\end{minipage} & \begin{minipage}[b]{\linewidth}\raggedright
Beitrag dieser Runde
\end{minipage} & \begin{minipage}[b]{\linewidth}\raggedright
Weiter offen
\end{minipage} \\
\midrule\noalign{}
\endhead
\bottomrule\noalign{}
\endlastfoot
T1 & Exakte Kommutanten-Leiter A→B→Cc→Cs→C; Clock als äußeres
Symmetrieelement & Welche Symmetrieoperationen der Compiler tatsächlich
bereitstellt \\
T2 & Lochbild bewiesen; Z\_h + Z\_add = 1; ε\_h, ε\_add & Konvergierter
Zustand, Renormierung, Grenzraum \\
T3 & - (ausgeklammert) & Räumlicher Träger \\
T4 & Kern-Symmetrieregel; kein einheitliches Vektor-Wörterbuch &
Chirales Maß, Anomalien \\
T5 & Exakte Krylov-Kette; ν₁\ldots ν₄ zweifach unabhängig &
Kontinuumslimes \\
T6 & Feldtyp-Census je Gradierung & Kopplungs- und Familiendaten \\
T7 & (1,0)⊕\allowbreak{}(0,1) ist B\_\{μν\}, kein Spin 2 & Dynamischer masseloser
Spin 2 \\
T8 & Sektorauswahl bewiesen bis g/Δ ≤ 0,0251; μ*-Fenster &
Zustandsfunktional, native Präparation \\
\end{longtable}
}

\textbf{Kein Tor wird als geschlossen markiert.} Keine RH-,
Faktorisierungs- oder P-versus-NP-Aussage; keine abgeleitete
Hylæan-Fähigkeit.

\subsection{8. Die jetzt entscheidenden
Prüfungen}\label{die-jetzt-entscheidenden-pruxfcfungen}

\begin{enumerate}
\def\labelenumi{\arabic{enumi}.}
\tightlist
\item
  \textbf{Untergrenzen verschärfen.} Die Casimir-Schranke trennt N=64
  bei 1/20 nicht von N=62 und N=66. Eine Jacobi-Vergleichsschranke mit
  rigorosen Normschranken für TT† je Bosonzahl oder eine Temple-Schranke
  auf der Ritz-Kette würde den Beweisbereich anheben.
\item
  \textbf{Den Zwei-Boson-Singulettraum ausschreiben.} w₂ ist der erste
  Zeuge; die volle Lanczos-Kette von H (nicht nur von T†) ab Tiefe 4
  braucht w₂ und seine Nachfolger.
\item
  \textbf{Operationen aus der Quelle.} Die Leiter sagt exakt, was
  so(10)-Operationen leisten würden; zu klären bleibt, welche davon der
  Compiler als markierte Operationen liefert. Der Clock ist ein äußeres
  Element und ersetzt sie nicht.
\item
  \textbf{Feldtyp vor Geometrie.} Jede chirale Aufspaltung macht den
  Bosontyp kanalabhängig (24/36 oder 40/20); eine räumliche Rechnung
  muss diese Aufteilung oder den einheitlichen Tensortyp explizit
  wählen.
\end{enumerate}

\subsection{9. Verifikation und nicht übernommene
Behauptungen}\label{verifikation-und-nicht-uxfcbernommene-behauptungen}

\codeword{replay.py} führt fünf Prüfer normal und unter \codeword{-OO}
aus und verlangt byteidentische Ausgaben; Status, Hashes und Guardsummen
stehen in \codeword{replay\_manifest.json}. Guardzahlen zählen
Prüfbedingungen, keine unabhängigen Theoreme. Exakt heißt
ganzzahlig/rational; numerisch heißt float64 mit Toleranz-Guard
(Ritz-Eigenwerte ab 3×3, Observablen, Erste-Moment-Auswertungen).

Nicht übernommen: der Worker-Beweis des N=64-Grundzustands bis 1/20;
eine physische Herleitung des Zustandsfunktionals; ein räumlicher
Träger; ein dynamischer Spin 2; die Existenz der Symmetrieoperationen
als Compiler- Primitive. Während der Runde festgestellte und behobene
Fehler der zunächst gelieferten Prüfer (u. a. Zyklus-Slots im
Spurnetzwerk, Bosonengewicht, doppelte T-Wirkung bei Bosonwiederholung,
Teilchen-Loch-Vorzeichen, min-max-Sektorschranke, skalierte w₂-Norm,
ungültiger Zykluszeuge) sind durch bleibende Guards abgesichert; keine
mathematische Aussage wurde an eine Erwartung angepasst.

Das \codeword{replay\_manifest.json} der Vorrunde v1.6.2 war durch einen
Lauf ohne Zugriff auf den externen Tensorpfad auf FAIL gesetzt worden;
es wurde aus den intakten \codeword{optimized.json} (PASS-Bedingung war
Byteidentität) rekonstruiert, Summen 27143/27139/4 stimmen mit dem
Delivery-Manifest.

\end{document}
